\section{The hard instances satisfy the containment conditions}
\label{sec:advantage}

Lowe et al.\ place normalized harmonic persistence in \(\BQP\) under two
access conditions~\cite[Lemma~7]{lowe2026normalized}: a polynomial-size
circuit preparing the uniform mixture over \(\cH_1=\ker\Delta_{1,d}\), and
the \emph{large overlap condition}, which asks that every nonzero eigenvalue
of \(P_{\cH_2}P_{\cH_1}P_{\cH_2}\) be at least \(a_{\min}\ge1/\poly\)
(their Definition~3).  This section shows that both conditions hold on the
instances produced by \cref{thm:general-hardness}, with \(a_{\min}\) close to
one, so that on this family the problem is simultaneously hard and
solvable in quantum polynomial time.  The mechanism is that the transfer
estimate, applied at zero energy, forces every harmonic vector to lie close
to the embedded logical kernel.

\subsection{Harmonic spaces track the logical kernels}

\begin{lemma}[Harmonic concentration]
\label{lem:harmonic-concentration}
Under the hypotheses and parameter choice \cref{eq:lambda-E-choice} of
\cref{thm:transfer}, every unit \(x\in\cH_A=\ker\Delta_A\) satisfies
\(\|(I-P_{K_A})x\|\le\eta_1\), where \(\eta_1=\sqrt C\,c_1\eta\le\eta/2\) after
fixing \(c_1\) with \(Cc_1^2\le1/4\).
\end{lemma}

\begin{proof}
\Cref{eq:transfer-bound} with \(\langle x,\Delta_Ax\rangle=0\) gives
\(\|(I-P_{K_A})x\|^2\le Ct_A\lambda^2\le Cc_1^2\eta^2\).
\end{proof}

\begin{lemma}[Equal dimensions]
\label{lem:equal-dimensions}
Let \(\cH,K\) be subspaces with \(\dim\cH=\dim K<\infty\), and suppose
\(\|(I-P_K)x\|\le s<1\) for every unit \(x\in\cH\).  Then every singular
value of \(P_K|_{\cH}\) is at least \(\sqrt{1-s^2}\), and
\(\|P_\cH-P_K\|\le s\).
\end{lemma}

\begin{proof}
For unit \(x\in\cH\), \(\|P_Kx\|^2=1-\|(I-P_K)x\|^2\ge1-s^2\), so
\(P_K|_\cH\) is injective and, by dimension count, maps \(\cH\) onto
\(K\).  The principal angles \(\theta_i\) between \(\cH\) and \(K\)
therefore satisfy \(\cos\theta_i\ge\sqrt{1-s^2}\), that is,
\(\sin\theta_i\le s\), and \(\|P_\cH-P_K\|=\sin\theta_{\max}\) for
subspaces of equal dimension.
\end{proof}

Since \(\dim\cH_A=\dim K_A\) by \cref{eq:kernel-gap-conclusion}, the two
lemmas give
\begin{equation}
\label{eq:projector-closeness}
\|P_{\cH_A}-P_{K_A}\|\le\eta_1\qquad(A\in\mathcal A).
\end{equation}
In words: the geometric harmonic space is an \(\eta_1\)-rotation of the
embedded logical kernel.  This is the approximate form of the isometry
requirement in Conjecture~2 of Lowe et al.; the exact rank statement is
\cref{thm:quotient}.

\subsection{Large overlap}

\begin{theorem}[Large overlap]
\label{thm:large-overlap}
For \(A\subseteq B\) in \(\mathcal A\), the operator
\(P_{\cH_B}P_{\cH_A}P_{\cH_B}\) has exactly \(\dim K_B\) nonzero eigenvalues,
and each of them is at least \((1-2\eta_1)^2\).  Hence the pair
\((P_{\cH_B},\cH_A)\) satisfies the large overlap condition with
\(a_{\min}\ge(1-2\eta_1)^2\ge1-2\eta\).
\end{theorem}

\begin{proof}
Because \(K_B\subseteq K_A\), \(P_{K_B}P_{K_A}=P_{K_B}\).  By
\cref{eq:projector-closeness},
\[
\|P_{\cH_B}P_{\cH_A}-P_{K_B}\|
\le\|P_{\cH_B}-P_{K_B}\|\,\|P_{\cH_A}\|+\|P_{K_B}\|\,\|P_{\cH_A}-P_{K_A}\|
\le2\eta_1 .
\]
By Weyl's inequality for singular values, every singular value of
\(P_{\cH_B}P_{\cH_A}\) is within \(2\eta_1\) of the corresponding singular
value of \(P_{K_B}\), which is \(1\) with multiplicity \(\dim K_B\) and \(0\)
otherwise.  The rank of \(P_{\cH_B}|_{\cH_A}\) is the rank of the induced
map \(H_d(X_A)\to H_d(X_B)\) (\cref{sec:prelim}), which equals \(\dim K_B\)
by \cref{thm:quotient}; so exactly \(\dim K_B\) singular values are nonzero,
each at least \(1-2\eta_1\).  The nonzero eigenvalues of
\(P_{\cH_B}P_{\cH_A}P_{\cH_B}\) are the squares of the nonzero singular values
of \(P_{\cH_B}P_{\cH_A}\).
\end{proof}

With \(\eta_0=1/10\) as in \cref{thm:general-hardness}, \(a_{\min}\ge4/5\);
choosing \(\eta=1/\poly\) instead gives \(a_{\min}\ge1-1/\poly\) at the price
of a smaller, still inverse-polynomial, gap \cref{eq:weighted-gap}.

\subsection{Preparing the initial harmonic mixture}

Let \(\rho_{K_{\rm in}}\) and \(\rho_{\cH_{\rm in}}\) denote the uniform
mixtures (normalized projectors) over \(K_{\rm in}=\cV C\subseteq C_d(\Xin)\)
and over \(\cH_{\rm in}=\ker\Delta_{\rm in}\).

\begin{theorem}[Preparation]
\label{thm:preparation}
The reduction of \cref{thm:general-hardness} can additionally output a
polynomial-size quantum circuit that prepares \(\rho_{K_{\rm in}}\) on the
simplex-label Hilbert space of \(\Xin\), up to \(1/\poly\) synthesis error.
Projecting its output onto \(\cH_{\rm in}\) by phase estimation of
\(\Delta_{\rm in}\) at resolution below \(E_{\rm in}\) succeeds with
probability at least \(1-\eta_1^2\) and returns a state \(\widetilde\rho\)
with
\[
\|\widetilde\rho-\rho_{\cH_{\rm in}}\|_1\le\frac{2\eta_1^2}{1-\eta_1^2}+\frac1{\poly}.
\]
\end{theorem}

\begin{proof}
The circuit has three stages.  First, prepare the fixed clean string, the
maximally mixed state on the \(r\) unconstrained qubits, and the clock state
\(Z^{-1/2}\sum_ts_t\ket t\) in the unary encoding; the amplitudes take the
two values \(Z^{-1/2}\) and \(\sqrt2Z^{-1/2}\), so this is a fixed real state
on \(L-1\) qubits with a polynomial-size description.  Second, for
\(t=1,\dots,L-1\) apply \(U_t\otimes\Pi^c_{\ge t}+I\otimes(I-\Pi^c_{\ge t})\),
where \(U_t\) is the unitary of step \(t\) (a gain or loss edge applies the
unitary \(H\); the scalars \(\sqrt2\) and \(1/\sqrt2\) are already in
\(s_t\)) and \(\Pi^c_{\ge t}\) projects onto clock values at least \(t\).
Exactly as in \cite[Lemma~8]{lowe2026normalized}, the clock branch \(\ket t\)
accumulates \(V_t=U_t\cdots U_1\), and the output is
\(\cV\rho_C\cV^*\) with \(\rho_C\) the uniform mixture over \(C\).  Third,
apply to every register qubit (clock and work) the fixed isometry
\(\C^2\to\C^{8}\) that sends the basis state \(\ket b\) to the normalized
oriented petal cycle \(\gamma_b\) of its bowtie, written on the labels of the
eight edges of that bowtie; the tensor product of these isometries is the
computational-basis embedding \((\C^2)^{\otimes n}\cong V\subseteq C_d(R)\subseteq C_d(\Xin)\)
of \cref{sec:prelim}, up to the orientation signs of join simplices, which
are computed classically from the vertex order and applied as phases.  The
result is \(\rho_{K_{\rm in}}\).

For the projection, \(\mathrm{Tr}(P_{\cH_{\rm in}}\rho_{K_{\rm in}})
=\frac1D\sum_{i=1}^D\cos^2\theta_i\ge1-\eta_1^2\) by
\cref{lem:equal-dimensions}, where \(\theta_i\) are the principal angles
between \(\cH_{\rm in}\) and \(K_{\rm in}\).  The post-measurement state is
\(\sum_i\cos^2\theta_i\,\ket{h_i}\!\bra{h_i}/\sum_j\cos^2\theta_j\) in the
orthonormal eigenbasis \((h_i)\) of \(P_{\cH}P_KP_{\cH}\) on \(\cH_{\rm in}\);
its eigenvalues lie in \([(1-\eta_1^2)/D,\,1/((1-\eta_1^2)D)]\), so its trace
distance to \(\rho_{\cH_{\rm in}}\) is at most \(2\eta_1^2/(1-\eta_1^2)\).
The Laplacian \(\Delta_{\rm in}\) is an explicit sparse operator on the
simplex labels with efficiently computable entries, and its phase estimation
at inverse-polynomial resolution is the standard primitive of quantum
topological data analysis~\cite{lloyd2016quantum,hayakawa2022quantum,mcardle2026streamlined};
the \(1/\poly\) term collects these approximation errors.
\end{proof}

\subsection{Consequences}

\begin{theorem}[Hard and in \(\BQP\)]
\label{thm:advantage}
Weighted gapped normalized persistence remains \(\SF^{G_2}\)-hard, and hence
\(\BQPone^{G_2}\)-hard and \(\SDQCone^{G_2}[\mathrm{frac}]\)-hard, when
restricted to instances satisfying
\begin{enumerate}[(a)]
\item the large overlap condition for \((P_{\cH_2},\cH_1)\) with
\(a_{\min}\ge1-1/\poly\), and
\item the existence of a polynomial-size circuit, computable in polynomial
time from the source instance, preparing a state within trace distance
\(1/\poly\) of the uniform mixture over \(\cH_1\).
\end{enumerate}
On instances satisfying (a) and (b) the problem is in \(\BQP\)
by \cite[Lemma~7]{lowe2026normalized}, whose proof is linear in the prepared
state and therefore tolerates the preparation error in (b).
\end{theorem}

\begin{proof}
\Cref{thm:general-hardness} with \(\eta=1/\poly\), \cref{thm:large-overlap},
and \cref{thm:preparation}.
\end{proof}

This is the format in which Lowe et al.\ state their own hardness results
(``hard even when the uniform mixture can be prepared and the large
overlap condition holds'', their Theorems~5, 7 and~8), and it is the standard
evidence for an exponential quantum advantage in this
literature~\cite{gyurik2022towards,gyurik2026provable}.  The unweighted
version is given in \cref{sec:unweighting}.
